\documentclass[a4paper,10pt]{article}
\usepackage{fontspec}
\usepackage{xeCJK}
\usepackage[a4paper,margin=20mm]{geometry}
\usepackage{graphicx}
\usepackage{fancyhdr}
\setsansfont[Path=/Users/yupyom/.src2tex/fonts/hackgen/, BoldFont=HackGen-Bold.ttf]{HackGen-Regular.ttf}
\setmonofont[Path=/Users/yupyom/.src2tex/fonts/hackgen/, BoldFont=HackGen-Bold.ttf]{HackGen-Regular.ttf}
\setCJKmainfont[Path=/usr/local/texlive/2026/texmf-dist/fonts/opentype/public/haranoaji/, Extension=.otf, BoldFont=HaranoAjiMincho-Bold.otf]{HaranoAjiMincho-Regular}
\setCJKsansfont[Path=/Users/yupyom/.src2tex/fonts/hackgen/, BoldFont=HackGen-Bold.ttf]{HackGen-Regular.ttf}
\setCJKmonofont[Path=/Users/yupyom/.src2tex/fonts/hackgen/, BoldFont=HackGen-Bold.ttf]{HackGen-Regular.ttf}
\XeTeXlinebreaklocale ""
\xeCJKsetup{CJKglue={},CJKecglue={}}
\newdimen\charwd
{\tt\global\setbox0=\hbox{x}\global\charwd=\wd0}
\frenchspacing
\pagestyle{fancy}
\renewcommand{\headrulewidth}{0pt}
\fancyhf{}
\fancyhead[R]{\rm src2tex version 2.236}
\fancyfoot[R]{\rm hanoi.php\qquad page \thepage}
% plain TeX compatibility
\makeatletter
\providecommand{\eqalign}[1]{\vcenter{\openup1\jot\m@th
  \ialign{\strut\hfil$\displaystyle{##}$&$\displaystyle{{}##}$\hfil\crcr#1\crcr}}}
\makeatother
\ifx\sevenrm\undefined
  \font\sevenrm=cmr7 scaled \magstep0
\fi
\def\mc{\relax}
\def\gt{\relax}
\def\sc{\scshape}
\begin{document}
\tt\mc 

\noindent
\mbox{}{\tt <}?php

\noindent
\mbox{}\rm\mc {\tt /}{\tt *}\  \hrulefill *
\tt\mc 

\noindent
\mbox{}\hfill

\noindent
\mbox{}\rm\mc \ % beginning of TeX mode

\centerline{\bf Towers of Hanoi (PHP)}

\begin{quote}

This program gives an answer to the following famous problem (towers of
Hanoi).
There is a legend that when one of the temples in Hanoi was constructed,
three poles were erected and a tower consisting of 64 golden discs was
arranged on one pole, their sizes decreasing regularly from bottom to top.
The monks were to move the tower of discs to the opposite pole, moving
only one at a time, and never putting any size disc above a smaller one.
The job was to be done in the minimum numbers of moves. What strategy for
moving discs will accomplish this optimum transfer?
\end{quote}

% end of TeX mode 
\tt\mc 

\noindent
\mbox{}\hfill

\noindent
\mbox{}\rm\mc {\tt *} \hrulefill \rm\mc \ {\tt *}{\tt /}
\tt\mc 

\noindent
\mbox{}\hfill

\noindent
\mbox{}\rm\mc {\tt /}{\tt *}\  \hrulefill\ hanoi.php \ \hrulefill \rm\mc \ {\tt *}{\tt /}
\tt\mc 

\noindent
\mbox{}\hfill

\noindent
\mbox{}{\tt\$}ARRAY{\tt\mc \ }={\tt\mc \ }8;{\tt\mc \ }{\tt\mc \ }{\tt\mc \ }{\tt\mc \ }{\tt\mc \ }{\tt\mc \ }{\tt\mc \ }{\tt\mc \ }{\tt\mc \ }{\tt\mc \ }{\tt\mc \ }{\tt\mc \ }{\tt\mc \ }{\tt\mc \ }{\tt\mc \ }{\tt\mc \ }{\tt\mc \ }{\tt\mc \ }{\tt\mc \ }{\tt\mc \ }{\tt\mc \ }{\tt\mc \ }{\tt\mc \ }{\tt\mc \ }{\tt\mc \ }{\tt\mc \ }{\tt\mc \ }{\tt\mc \ }{\tt\mc \ }\rm\mc {\tt /}{\tt *}\  \ disc の数 \hfill \rm\mc \ {\tt *}{\tt /}
\tt\mc 

\noindent
\mbox{}\hfill

\noindent
\mbox{}{\tt\$}disc{\tt\mc \ }={\tt\mc \ }array{\tt\_}fill(0,{\tt\mc \ }3,{\tt\mc \ }array{\tt\_}fill(0,{\tt\mc \ }{\tt\$}ARRAY,{\tt\mc \ }0));

\noindent
\mbox{}{\tt\mc \ }{\tt\mc \ }{\tt\mc \ }{\tt\mc \ }{\tt\mc \ }{\tt\mc \ }{\tt\mc \ }{\tt\mc \ }{\tt\mc \ }{\tt\mc \ }{\tt\mc \ }{\tt\mc \ }{\tt\mc \ }{\tt\mc \ }{\tt\mc \ }{\tt\mc \ }{\tt\mc \ }{\tt\mc \ }{\tt\mc \ }{\tt\mc \ }{\tt\mc \ }{\tt\mc \ }{\tt\mc \ }{\tt\mc \ }{\tt\mc \ }{\tt\mc \ }{\tt\mc \ }{\tt\mc \ }{\tt\mc \ }{\tt\mc \ }{\tt\mc \ }{\tt\mc \ }{\tt\mc \ }{\tt\mc \ }{\tt\mc \ }{\tt\mc \ }{\tt\mc \ }{\tt\mc \ }{\tt\mc \ }{\tt\mc \ }\rm\mc {\tt /}{\tt *}\  \ disc に関するデータの置き場所\hfill \rm\mc \ {\tt *}{\tt /}
\tt\mc 

\noindent
\mbox{}\hfill

\noindent
\mbox{}\textbf{function}{\tt\mc \ }initArray(){\tt\mc \ }{\tt\char'173}{\tt\mc \ }{\tt\mc \ }{\tt\mc \ }{\tt\mc \ }{\tt\mc \ }{\tt\mc \ }{\tt\mc \ }{\tt\mc \ }{\tt\mc \ }{\tt\mc \ }{\tt\mc \ }{\tt\mc \ }{\tt\mc \ }{\tt\mc \ }{\tt\mc \ }{\tt\mc \ }{\tt\mc \ }{\tt\mc \ }\rm\mc {\tt /}{\tt *}\  \ disc に関するデータの初期化\hfill \rm\mc \ {\tt *}{\tt /}
\tt\mc 

\noindent
\mbox{}{\tt\mc \ }{\tt\mc \ }{\tt\mc \ }{\tt\mc \ }global{\tt\mc \ }{\tt\$}ARRAY,{\tt\mc \ }{\tt\$}disc;

\noindent
\mbox{}{\tt\mc \ }{\tt\mc \ }{\tt\mc \ }{\tt\mc \ }\textbf{for}{\tt\mc \ }({\tt\$}j{\tt\mc \ }={\tt\mc \ }0;{\tt\mc \ }{\tt\$}j{\tt\mc \ }{\tt <}{\tt\mc \ }{\tt\$}ARRAY;{\tt\mc \ }{\tt\$}j++){\tt\mc \ }{\tt\char'173}

\noindent
\mbox{}{\tt\mc \ }{\tt\mc \ }{\tt\mc \ }{\tt\mc \ }{\tt\mc \ }{\tt\mc \ }{\tt\mc \ }{\tt\mc \ }{\tt\$}disc[0][{\tt\$}j]{\tt\mc \ }={\tt\mc \ }{\tt\$}ARRAY{\tt\mc \ }{\tt -}{\tt\mc \ }{\tt\$}j;

\noindent
\mbox{}{\tt\mc \ }{\tt\mc \ }{\tt\mc \ }{\tt\mc \ }{\tt\mc \ }{\tt\mc \ }{\tt\mc \ }{\tt\mc \ }{\tt\$}disc[1][{\tt\$}j]{\tt\mc \ }={\tt\mc \ }0;

\noindent
\mbox{}{\tt\mc \ }{\tt\mc \ }{\tt\mc \ }{\tt\mc \ }{\tt\mc \ }{\tt\mc \ }{\tt\mc \ }{\tt\mc \ }{\tt\$}disc[2][{\tt\$}j]{\tt\mc \ }={\tt\mc \ }0;

\noindent
\mbox{}{\tt\mc \ }{\tt\mc \ }{\tt\mc \ }{\tt\mc \ }{\tt\char'175}

\noindent
\mbox{}{\tt\char'175}

\noindent
\mbox{}\hfill

\noindent
\mbox{}{\tt\$}counter{\tt\mc \ }={\tt\mc \ }0;{\tt\mc \ }{\tt\mc \ }{\tt\mc \ }{\tt\mc \ }{\tt\mc \ }{\tt\mc \ }{\tt\mc \ }{\tt\mc \ }{\tt\mc \ }{\tt\mc \ }{\tt\mc \ }{\tt\mc \ }{\tt\mc \ }{\tt\mc \ }{\tt\mc \ }{\tt\mc \ }{\tt\mc \ }{\tt\mc \ }{\tt\mc \ }{\tt\mc \ }{\tt\mc \ }{\tt\mc \ }{\tt\mc \ }{\tt\mc \ }{\tt\mc \ }{\tt\mc \ }{\tt\mc \ }\rm\mc {\tt /}{\tt *}\  \ 移動回数カウンタ \hfill \rm\mc \ {\tt *}{\tt /}
\tt\mc 

\noindent
\mbox{}\hfill

\noindent
\mbox{}\textbf{function}{\tt\mc \ }printResult(){\tt\mc \ }{\tt\char'173}{\tt\mc \ }{\tt\mc \ }{\tt\mc \ }{\tt\mc \ }{\tt\mc \ }{\tt\mc \ }{\tt\mc \ }{\tt\mc \ }{\tt\mc \ }{\tt\mc \ }{\tt\mc \ }{\tt\mc \ }{\tt\mc \ }{\tt\mc \ }{\tt\mc \ }{\tt\mc \ }\rm\mc {\tt /}{\tt *}\  \ 結果の表示 \hfill \rm\mc \ {\tt *}{\tt /}
\tt\mc 

\noindent
\mbox{}{\tt\mc \ }{\tt\mc \ }{\tt\mc \ }{\tt\mc \ }global{\tt\mc \ }{\tt\$}ARRAY,{\tt\mc \ }{\tt\$}disc,{\tt\mc \ }{\tt\$}counter;

\noindent
\mbox{}{\tt\mc \ }{\tt\mc \ }{\tt\mc \ }{\tt\mc \ }{\tt\$}counter++;

\noindent
\mbox{}{\tt\mc \ }{\tt\mc \ }{\tt\mc \ }{\tt\mc \ }printf({\tt "}{\tt -}{\tt -}{\tt -}{\tt\#}{\tt\%}d{\tt -}{\tt -}{\tt -}{\tt\char92}n{\tt "},{\tt\mc \ }{\tt\$}counter);

\noindent
\mbox{}{\tt\mc \ }{\tt\mc \ }{\tt\mc \ }{\tt\mc \ }\textbf{for}{\tt\mc \ }({\tt\$}i{\tt\mc \ }={\tt\mc \ }0;{\tt\mc \ }{\tt\$}i{\tt\mc \ }{\tt <}={\tt\mc \ }2;{\tt\mc \ }{\tt\$}i++){\tt\mc \ }{\tt\char'173}

\noindent
\mbox{}{\tt\mc \ }{\tt\mc \ }{\tt\mc \ }{\tt\mc \ }{\tt\mc \ }{\tt\mc \ }{\tt\mc \ }{\tt\mc \ }printf({\tt "}[{\tt\%}d]{\tt\mc \ }{\tt "},{\tt\mc \ }{\tt\$}i);

\noindent
\mbox{}{\tt\mc \ }{\tt\mc \ }{\tt\mc \ }{\tt\mc \ }{\tt\mc \ }{\tt\mc \ }{\tt\mc \ }{\tt\mc \ }\textbf{for}{\tt\mc \ }({\tt\$}j{\tt\mc \ }={\tt\mc \ }0;{\tt\mc \ }{\tt\$}j{\tt\mc \ }{\tt <}{\tt\mc \ }{\tt\$}ARRAY;{\tt\mc \ }{\tt\$}j++){\tt\mc \ }{\tt\char'173}

\noindent
\mbox{}{\tt\mc \ }{\tt\mc \ }{\tt\mc \ }{\tt\mc \ }{\tt\mc \ }{\tt\mc \ }{\tt\mc \ }{\tt\mc \ }{\tt\mc \ }{\tt\mc \ }{\tt\mc \ }{\tt\mc \ }\textbf{if}{\tt\mc \ }({\tt\$}disc[{\tt\$}i][{\tt\$}j]{\tt\mc \ }!={\tt\mc \ }0){\tt\mc \ }{\tt\char'173}

\noindent
\mbox{}{\tt\mc \ }{\tt\mc \ }{\tt\mc \ }{\tt\mc \ }{\tt\mc \ }{\tt\mc \ }{\tt\mc \ }{\tt\mc \ }{\tt\mc \ }{\tt\mc \ }{\tt\mc \ }{\tt\mc \ }{\tt\mc \ }{\tt\mc \ }{\tt\mc \ }{\tt\mc \ }printf({\tt "}{\tt\%}d{\tt\mc \ }{\tt "},{\tt\mc \ }{\tt\$}disc[{\tt\$}i][{\tt\$}j]);

\noindent
\mbox{}{\tt\mc \ }{\tt\mc \ }{\tt\mc \ }{\tt\mc \ }{\tt\mc \ }{\tt\mc \ }{\tt\mc \ }{\tt\mc \ }{\tt\mc \ }{\tt\mc \ }{\tt\mc \ }{\tt\mc \ }{\tt\char'175}{\tt\mc \ }\textbf{else}{\tt\mc \ }{\tt\char'173}

\noindent
\mbox{}{\tt\mc \ }{\tt\mc \ }{\tt\mc \ }{\tt\mc \ }{\tt\mc \ }{\tt\mc \ }{\tt\mc \ }{\tt\mc \ }{\tt\mc \ }{\tt\mc \ }{\tt\mc \ }{\tt\mc \ }{\tt\mc \ }{\tt\mc \ }{\tt\mc \ }{\tt\mc \ }\textbf{break};

\noindent
\mbox{}{\tt\mc \ }{\tt\mc \ }{\tt\mc \ }{\tt\mc \ }{\tt\mc \ }{\tt\mc \ }{\tt\mc \ }{\tt\mc \ }{\tt\mc \ }{\tt\mc \ }{\tt\mc \ }{\tt\mc \ }{\tt\char'175}

\noindent
\mbox{}{\tt\mc \ }{\tt\mc \ }{\tt\mc \ }{\tt\mc \ }{\tt\mc \ }{\tt\mc \ }{\tt\mc \ }{\tt\mc \ }{\tt\char'175}

\noindent
\mbox{}{\tt\mc \ }{\tt\mc \ }{\tt\mc \ }{\tt\mc \ }{\tt\mc \ }{\tt\mc \ }{\tt\mc \ }{\tt\mc \ }printf({\tt "}{\tt\char92}n{\tt "});

\noindent
\mbox{}{\tt\mc \ }{\tt\mc \ }{\tt\mc \ }{\tt\mc \ }{\tt\char'175}

\noindent
\mbox{}{\tt\char'175}

\noindent
\mbox{}\hfill

\noindent
\mbox{}{\tt\$}ptr{\tt\mc \ }={\tt\mc \ }[0,{\tt\mc \ }0,{\tt\mc \ }0];{\tt\mc \ }{\tt\mc \ }{\tt\mc \ }{\tt\mc \ }{\tt\mc \ }{\tt\mc \ }{\tt\mc \ }{\tt\mc \ }{\tt\mc \ }{\tt\mc \ }{\tt\mc \ }{\tt\mc \ }{\tt\mc \ }{\tt\mc \ }{\tt\mc \ }{\tt\mc \ }{\tt\mc \ }{\tt\mc \ }{\tt\mc \ }{\tt\mc \ }{\tt\mc \ }{\tt\mc \ }{\tt\mc \ }\rm\mc {\tt /}{\tt *}\  \ disc 移動用ポインタ（インデックス）\hfill \rm\mc \ {\tt *}{\tt /}
\tt\mc 

\noindent
\mbox{}\hfill

\noindent
\mbox{}\textbf{function}{\tt\mc \ }moveOneDisc({\tt\$}i,{\tt\mc \ }{\tt\$}j){\tt\mc \ }{\tt\char'173}{\tt\mc \ }{\tt\mc \ }{\tt\mc \ }{\tt\mc \ }{\tt\mc \ }{\tt\mc \ }{\tt\mc \ }{\tt\mc \ }{\tt\mc \ }{\tt\mc \ }\rm\mc {\tt /}{\tt *}\  \ 1枚の disc を pole $i$ からpole $j$ に移動する \hfill \rm\mc \ {\tt *}{\tt /}
\tt\mc 

\noindent
\mbox{}{\tt\mc \ }{\tt\mc \ }{\tt\mc \ }{\tt\mc \ }global{\tt\mc \ }{\tt\$}disc,{\tt\mc \ }{\tt\$}ptr;

\noindent
\mbox{}{\tt\mc \ }{\tt\mc \ }{\tt\mc \ }{\tt\mc \ }{\tt\$}ptr[{\tt\$}i]{\tt -}{\tt -};

\noindent
\mbox{}{\tt\mc \ }{\tt\mc \ }{\tt\mc \ }{\tt\mc \ }{\tt\$}disc[{\tt\$}j][{\tt\$}ptr[{\tt\$}j]]{\tt\mc \ }={\tt\mc \ }{\tt\$}disc[{\tt\$}i][{\tt\$}ptr[{\tt\$}i]];

\noindent
\mbox{}{\tt\mc \ }{\tt\mc \ }{\tt\mc \ }{\tt\mc \ }{\tt\$}ptr[{\tt\$}j]++;

\noindent
\mbox{}{\tt\mc \ }{\tt\mc \ }{\tt\mc \ }{\tt\mc \ }{\tt\$}disc[{\tt\$}i][{\tt\$}ptr[{\tt\$}i]]{\tt\mc \ }={\tt\mc \ }0;

\noindent
\mbox{}{\tt\char'175}

\noindent
\mbox{}\hfill

\noindent
\mbox{}\textbf{function}{\tt\mc \ }moveDiscs({\tt\$}n,{\tt\mc \ }{\tt\$}i,{\tt\mc \ }{\tt\$}j,{\tt\mc \ }{\tt\$}k){\tt\mc \ }{\tt\char'173}{\tt\mc \ }{\tt\mc \ }{\tt\mc \ }{\tt\mc \ }\rm\mc {\tt /}{\tt *}\  \ \underline{\textsf{上から $n$ 枚目までの disc}}を、pole $i$ から pole $j$ に\hfill \rm\mc \ {\tt *}{\tt /}
\tt\mc 

\noindent
\mbox{}{\tt\mc \ }{\tt\mc \ }{\tt\mc \ }{\tt\mc \ }{\tt\mc \ }{\tt\mc \ }{\tt\mc \ }{\tt\mc \ }{\tt\mc \ }{\tt\mc \ }{\tt\mc \ }{\tt\mc \ }{\tt\mc \ }{\tt\mc \ }{\tt\mc \ }{\tt\mc \ }{\tt\mc \ }{\tt\mc \ }{\tt\mc \ }{\tt\mc \ }{\tt\mc \ }{\tt\mc \ }{\tt\mc \ }{\tt\mc \ }{\tt\mc \ }{\tt\mc \ }{\tt\mc \ }{\tt\mc \ }{\tt\mc \ }{\tt\mc \ }{\tt\mc \ }{\tt\mc \ }{\tt\mc \ }{\tt\mc \ }{\tt\mc \ }{\tt\mc \ }{\tt\mc \ }{\tt\mc \ }{\tt\mc \ }{\tt\mc \ }\rm\mc {\tt /}{\tt *}\  \ pole $k$ を経由して、移動する\hfill \rm\mc \ {\tt *}{\tt /}
\tt\mc 

\noindent
\mbox{}{\tt\mc \ }{\tt\mc \ }{\tt\mc \ }{\tt\mc \ }\textbf{if}{\tt\mc \ }({\tt\$}n{\tt\mc \ }{\tt >}={\tt\mc \ }1){\tt\mc \ }{\tt\char'173}

\noindent
\mbox{}{\tt\mc \ }{\tt\mc \ }{\tt\mc \ }{\tt\mc \ }{\tt\mc \ }{\tt\mc \ }{\tt\mc \ }{\tt\mc \ }moveDiscs({\tt\$}n{\tt\mc \ }{\tt -}{\tt\mc \ }1,{\tt\mc \ }{\tt\$}i,{\tt\mc \ }{\tt\$}k,{\tt\mc \ }{\tt\$}j);{\tt\mc \ }{\tt\mc \ }\rm\mc {\tt /}{\tt *}\  \ 関数 {\tt moveDiscs()}の中で、さらに自分自身 \hfill \rm\mc \ {\tt *}{\tt /}
\tt\mc 

\noindent
\mbox{}{\tt\mc \ }{\tt\mc \ }{\tt\mc \ }{\tt\mc \ }{\tt\mc \ }{\tt\mc \ }{\tt\mc \ }{\tt\mc \ }moveOneDisc({\tt\$}i,{\tt\mc \ }{\tt\$}j);{\tt\mc \ }{\tt\mc \ }{\tt\mc \ }{\tt\mc \ }{\tt\mc \ }{\tt\mc \ }{\tt\mc \ }{\tt\mc \ }{\tt\mc \ }{\tt\mc \ }{\tt\mc \ }{\tt\mc \ }\rm\mc {\tt /}{\tt *}\  \ {\tt moveDiscs()} が使われている。このような \hfill \rm\mc \ {\tt *}{\tt /}
\tt\mc 

\noindent
\mbox{}{\tt\mc \ }{\tt\mc \ }{\tt\mc \ }{\tt\mc \ }{\tt\mc \ }{\tt\mc \ }{\tt\mc \ }{\tt\mc \ }printResult();{\tt\mc \ }{\tt\mc \ }{\tt\mc \ }{\tt\mc \ }{\tt\mc \ }{\tt\mc \ }{\tt\mc \ }{\tt\mc \ }{\tt\mc \ }{\tt\mc \ }{\tt\mc \ }{\tt\mc \ }{\tt\mc \ }{\tt\mc \ }{\tt\mc \ }{\tt\mc \ }{\tt\mc \ }{\tt\mc \ }\rm\mc {\tt /}{\tt *}\  \ 手法は、「再帰的呼びだし」といわれる。 \hfill \rm\mc \ {\tt *}{\tt /}
\tt\mc 

\noindent
\mbox{}{\tt\mc \ }{\tt\mc \ }{\tt\mc \ }{\tt\mc \ }{\tt\mc \ }{\tt\mc \ }{\tt\mc \ }{\tt\mc \ }moveDiscs({\tt\$}n{\tt\mc \ }{\tt -}{\tt\mc \ }1,{\tt\mc \ }{\tt\$}k,{\tt\mc \ }{\tt\$}j,{\tt\mc \ }{\tt\$}i);

\noindent
\mbox{}{\tt\mc \ }{\tt\mc \ }{\tt\mc \ }{\tt\mc \ }{\tt\char'175}

\noindent
\mbox{}{\tt\char'175}

\noindent
\mbox{}\hfill

\noindent
\mbox{}\rm\mc {\tt /}{\tt *}\  \par\begin{center}

\includegraphics[scale=0.3]{hanoi1}\quad
\includegraphics[scale=0.3]{hanoi2}\end{center}

たとえば、関数 {\tt moveDiscs(4, 0, 1, 2)} を呼び出すことは、
上図のような操作をすることに対応する。\hfill \rm\mc \ {\tt *}{\tt /}
\tt\mc 

\noindent
\mbox{}\hfill

\noindent
\mbox{}{\tt\$}ptr[0]{\tt\mc \ }={\tt\mc \ }{\tt\$}ARRAY;

\noindent
\mbox{}{\tt\$}ptr[1]{\tt\mc \ }={\tt\mc \ }0;

\noindent
\mbox{}{\tt\$}ptr[2]{\tt\mc \ }={\tt\mc \ }0;

\noindent
\mbox{}\hfill

\noindent
\mbox{}initArray();

\noindent
\mbox{}moveDiscs({\tt\$}ARRAY,{\tt\mc \ }0,{\tt\mc \ }1,{\tt\mc \ }2);{\tt\mc \ }{\tt\mc \ }{\tt\mc \ }{\tt\mc \ }{\tt\mc \ }{\tt\mc \ }{\tt\mc \ }{\tt\mc \ }{\tt\mc \ }{\tt\mc \ }{\tt\mc \ }{\tt\mc \ }{\tt\mc \ }\rm\mc {\tt /}{\tt *}\  \ {\tt ARRAY} 枚の disc をpole 0 から pole 1 に pole 2\hfill \rm\mc \ {\tt *}{\tt /}
\tt\mc 

\noindent
\mbox{}{\tt\mc \ }{\tt\mc \ }{\tt\mc \ }{\tt\mc \ }{\tt\mc \ }{\tt\mc \ }{\tt\mc \ }{\tt\mc \ }{\tt\mc \ }{\tt\mc \ }{\tt\mc \ }{\tt\mc \ }{\tt\mc \ }{\tt\mc \ }{\tt\mc \ }{\tt\mc \ }{\tt\mc \ }{\tt\mc \ }{\tt\mc \ }{\tt\mc \ }{\tt\mc \ }{\tt\mc \ }{\tt\mc \ }{\tt\mc \ }{\tt\mc \ }{\tt\mc \ }{\tt\mc \ }{\tt\mc \ }{\tt\mc \ }{\tt\mc \ }{\tt\mc \ }{\tt\mc \ }{\tt\mc \ }{\tt\mc \ }{\tt\mc \ }{\tt\mc \ }{\tt\mc \ }{\tt\mc \ }{\tt\mc \ }{\tt\mc \ }\rm\mc {\tt /}{\tt *}\  \ を経由して、移動する \hfill \rm\mc \ {\tt *}{\tt /}

\tt\mc 

\rm\mc

\end{document}
